% Corrected French for printed p.176; departures are append-only below.
\SourcePage{176}

\LayerTarget{tohoku:corollary:3.8.1:1:p2}
\noindent\textit{tous les $U_{i_0\ldots i_p}$ sont $F$-acycliques
\textup{(}c'est-à-dire satisfont $H^q(U_{i_0\ldots i_p},F)=0$ pour $q>0$\textup{)}.}

\LayerTarget{tohoku:def:point-indexed-cover-system}
Bornons-nous maintenant aux recouvrements
$\mathbf U=(U_x)_{x\in X}$ indexés par les points $x$ de $X$, tels que
$x\in U_x$ pour tout $x$, ordonnons-les en écrivant
$\mathbf U\leqq\mathbf U'$ si $U'_x\subset U_x$ pour tout $x$.
\LayerTarget{tohoku:map:point-cover-refinement-complexes}
Si $C_{\mathrm U}$ et $C_{\mathrm U'}$ sont les foncteurs-complexes
correspondants sur $\mathbf C^X$, on aura alors un homomorphisme fonctoriel
$C_{\mathrm U}\to C_{\mathrm U'}$, d'où un homomorphisme pour
les foncteurs spectraux correspondants. Un passage immédiat à la limite
inductive donne alors le

\LayerTarget{tohoku:corollary:3.8.1:2}
\textsc{Corollaire 2.} \textit{Soit $X$ un espace topologique
quelconque. Il existe un foncteur spectral sur la catégorie $\mathbf C^X$ des faisceaux
de groupes abéliens sur $X$, aboutissant au foncteur gradué $(H^n(X,F))$, et
dont le terme initial est donné par}
\LayerTarget{tohoku:eq:global-cech-spectral-E2}
\[
E_2^{p,q}(F)=\check H^p(X,H^q(F))
\]
\textit{\textup{(}$H^q(F)$ étant le préfaisceau défini dans le th. \textup{3.8.1)}. On a}
\LayerTarget{tohoku:eq:global-cech-spectral-E20q}
\[
E_2^{0,q}(F)=0\quad\textit{pour }q>0.
\]

Cette dernière formule résulte de la définition de
$\check H^0(X,H^q(F))$ et du fait suivant :

\LayerTarget{tohoku:lemma:3.8.2}
\textsc{Lemme 3.8.2.} \textit{Soit $U$ un voisinage de $x$ et
soit $c^q\in H^q(U,F)$, avec $q>0$, alors il existe un voisinage $V\subset U$ de $x$ tel
que l'image de $c^q$ dans $H^q(V,F)$ soit nulle.}

\LayerTarget{tohoku:proof:3.8.2}
Pour le voir, il suffit de prendre une résolution injective de $F\mid U$, soit
$C$ le complexe correspondant, et de représenter $c^q$ par un élément de
$H^q(\Gamma(U,C))$, défini par un cocycle $z\in\Gamma(U,C^q)$ ; d'après
l'acyclicité de $C$ en dimension $q$, la restriction de $z$ à un voisinage
convenable $V$ de $x$ est un cobord, d'où le résultat, en remarquant que
$C\mid V$ est une résolution injective de $F\mid V$ en vertu de prop. 3.1.3,
donc que $H^q(V,F)=H^q(\Gamma(C\mid V))$.

La suite spectrale du corollaire 2 donne des homomorphismes fonctoriels
\LayerTarget{tohoku:map:global-cech-to-derived-cohomology}
\[
\check H^p(X,F)\to H^p(X,F)
\]
et la formule
\LayerTarget{tohoku:eq:global-cech-E201}
$E_2^{0,1}=0$
montre que :

\LayerTarget{tohoku:corollary:3.8.1:3}
\textsc{Corollaire 3.} \textit{Les homomorphismes précédents sont
bijectifs si $p=0$ ou $1$ \textup{(}ce que nous savions déjà\textup{)} et des monomorphismes si
$p=2$.}

\LayerTarget{tohoku:remark:paracompact-cech-comparison}
On retrouve d'ailleurs le fait que \textit{si $X$ est paracompact}, alors
$\check H^p=H^p$, plus précisément, \textit{les homomorphismes canoniques
ci-dessus sont alors des isomorphismes.} En effet, on vérifie dans ce cas
(grâce au fait que le faisceau associé au préfaisceau $H^q(F)$ est nul si
$q>0$) que
\LayerTarget{tohoku:eq:paracompact-global-cech-E2-vanishing}
$E_2^{p,q}=0$ pour $q>0$,
cf. [9].

\LayerTarget{tohoku:criterion:global-cech-edge-range}
Le corollaire 3 se généralise ainsi : \textit{Si
$\check H^p(X,H^q(F))=0$ pour $0<q<n$, alors l'homomorphisme
$\check H^i(X,F)\to H^i(X,F)$ est un isomorphisme pour $i\leqq n$
et un monomorphisme pour $i=n+1$.} On en conclut (avec H. Cartan) :

\LayerTarget{tohoku:corollary:3.8.1:4:p1}
\textsc{Corollaire 4.} \textit{Soit $\mathbf U$ un ensemble
d'ouverts formant une base pour la topologie de $X$, soit $F$ un faisceau
abélien sur $X$ tel que pour toute suite non}

\TohokuContinuation{177}
